\documentclass[12pt,a4paper]{article}
\usepackage[utf8]{inputenc}
\usepackage[russian]{babel}
\usepackage{geometry}
\usepackage{listings}
\usepackage{xcolor}
\usepackage{hyperref}
\usepackage{graphicx}
\usepackage{amsmath}

\geometry{margin=2.5cm}

% Настройка подсветки синтаксиса
\lstset{
    language=C++,
    basicstyle=\ttfamily\small,
    keywordstyle=\color{blue}\bfseries,
    commentstyle=\color{green!60!black},
    stringstyle=\color{red},
    numberstyle=\tiny\color{gray},
    breaklines=true,
    frame=single,
    numbers=left,
    stepnumber=1,
    tabsize=4,
    showspaces=false,
    showstringspaces=false
}

\title{Руководство пользователя\\Компилятор языка my\_lang}
\author{Версия 1.0}
\date{\today}

\begin{document}

\maketitle
\tableofcontents
\newpage

\section{Введение}

Данное руководство описывает процесс использования компилятора языка программирования \texttt{my\_lang}. Компилятор выполняет лексический анализ, синтаксический разбор, семантическую проверку и генерацию промежуточного представления в виде обратной польской записи (POLIZ).

\section{Системные требования}

Для работы компилятора необходимы:
\begin{itemize}
    \item Операционная система: Windows, Linux или macOS
    \item Компилятор C++ с поддержкой стандарта C++20 (GCC, Clang, MSVC)
    \item CMake версии 3.20 или выше
    \item Python 3.x (для генерации парсера и лексера)
    \item GoogleTest (автоматически загружается через CMake)
\end{itemize}

\section{Установка}

\subsection{Получение исходного кода}

Исходный код проекта должен находиться в директории \texttt{c:\project} (или другой выбранной директории).

\subsection{Сборка проекта}

Для сборки проекта выполните следующие команды:

\begin{lstlisting}[language=bash]
cd c:\project
mkdir cmake-build-debug
cd cmake-build-debug
cmake ..
cmake --build .
\end{lstlisting}

При сборке автоматически выполняются следующие действия:
\begin{itemize}
    \item Генерация парсера из файла \texttt{grammar.txt}
    \item Генерация лексера из файла \texttt{system\_reserved\_identifiers.txt}
    \item Компиляция всех исходных файлов
    \item Сборка исполняемого файла \texttt{my\_lang\_ftl\_prak.exe}
\end{itemize}

\section{Использование компилятора}

\subsection{Базовое использование}

Для компиляции программы на языке \texttt{my\_lang} используйте следующую команду:

\begin{lstlisting}[language=bash]
./my_lang_ftl_prak.exe <путь_к_файлу>
\end{lstlisting}

Если путь к файлу не указан, по умолчанию используется файл \texttt{test\_program.txt}.

\subsection{Пример использования}

\begin{lstlisting}[language=bash]
./my_lang_ftl_prak.exe test_program.txt
\end{lstlisting}

\subsection{Выходные данные}

Компилятор выводит следующую информацию:
\begin{itemize}
    \item Путь к обрабатываемому файлу
    \item Результаты семантической проверки
    \item Результаты оптимизации AST
    \item Сгенерированный POLIZ код
\end{itemize}

При успешной компиляции вы увидите сообщения:
\begin{itemize}
    \item \texttt{✅ Semantic checks passed.}
    \item \texttt{✅ AST optimization complete.}
    \item \texttt{✅ POLIZ generation complete.}
\end{itemize}

При ошибках компиляции выводится соответствующее сообщение с описанием проблемы.

\section{Синтаксис языка my\_lang}

\subsection{Типы данных}

Язык поддерживает следующие базовые типы:
\begin{itemize}
    \item \texttt{int} --- целые числа
    \item \texttt{float} --- числа с плавающей точкой одинарной точности
    \item \texttt{double} --- числа с плавающей точкой двойной точности
    \item \texttt{string} --- строки
    \item \texttt{vector<T>} --- векторы (динамические массивы) элементов типа T
\end{itemize}

\subsection{Объявление переменных}

Переменные объявляются следующим образом:

\begin{lstlisting}
int x;
double y = 3.14;
string name = "Hello";
vector<int> numbers;
\end{lstlisting}

\subsection{Объявление функций}

Функции объявляются с указанием типа возвращаемого значения:

\begin{lstlisting}
int add(int a, int b) {
    return a + b;
}
\end{lstlisting}

\subsection{Объявление классов}

Классы объявляются с помощью ключевого слова \texttt{class}:

\begin{lstlisting}
class Person {
    int age;
    string name;
    
    void print() {
        print(name, age);
    }
};
\end{lstlisting}

\subsection{Управляющие конструкции}

\subsubsection{Условный оператор if}

\begin{lstlisting}
if (x > 0) {
    print("Positive");
} else {
    print("Non-positive");
}
\end{lstlisting}

\subsubsection{Цикл while}

\begin{lstlisting}
while (i < 10) {
    i = i + 1;
}
\end{lstlisting}

\subsubsection{Цикл for}

\begin{lstlisting}
for (i < 10) {
    i = i + 1;
}
\end{lstlisting}

\subsubsection{Оператор after}

Оператор \texttt{after} позволяет выполнить код после завершения цикла:

\begin{lstlisting}
after (i < 10) {
    print("Loop finished");
}
\end{lstlisting}

\subsection{Операторы}

\subsubsection{Арифметические операторы}

\begin{itemize}
    \item \texttt{+} --- сложение
    \item \texttt{-} --- вычитание
    \item \texttt{*} --- умножение
    \item \texttt{/} --- деление
    \item \texttt{\%} --- остаток от деления
\end{itemize}

\subsubsection{Операторы сравнения}

\begin{itemize}
    \item \texttt{==} --- равно
    \item \texttt{!=} --- не равно
    \item \texttt{<} --- меньше
    \item \texttt{>} --- больше
    \item \texttt{<=} --- меньше или равно
    \item \texttt{>=} --- больше или равно
\end{itemize}

\subsubsection{Логические операторы}

\begin{itemize}
    \item \texttt{\&\&} --- логическое И
    \item \texttt{||} --- логическое ИЛИ
    \item \texttt{!} --- логическое НЕ
    \item \texttt{or} --- альтернативная форма ИЛИ
\end{itemize}

\subsubsection{Операторы присваивания}

\begin{itemize}
    \item \texttt{=} --- присваивание
    \item \texttt{+=} --- сложение с присваиванием
    \item \texttt{-=} --- вычитание с присваиванием
    \item \texttt{*=} --- умножение с присваиванием
    \item \texttt{/=} --- деление с присваиванием
    \item \texttt{\%=} --- остаток с присваиванием
\end{itemize}

\subsubsection{Инкремент и декремент}

\begin{itemize}
    \item \texttt{++} --- инкремент (префиксный и постфиксный)
    \item \texttt{--} --- декремент (префиксный и постфиксный)
\end{itemize}

\subsection{Встроенные функции}

\subsubsection{print}

Функция \texttt{print} выводит значения на экран:

\begin{lstlisting}
print("Hello, World!");
print(x, y, z);
\end{lstlisting}

\subsubsection{input}

Функция \texttt{input} предназначена для ввода данных (в текущей реализации может быть не полностью реализована).

\subsection{Доступ к полям объектов}

Доступ к полям объектов осуществляется через точку:

\begin{lstlisting}
person.name = "Alice";
person.age = 25;
\end{lstlisting}

\subsection{Доступ к элементам массивов}

Доступ к элементам массивов и векторов осуществляется через квадратные скобки:

\begin{lstlisting}
vector<int> arr;
arr[0] = 10;
int value = arr[0];
\end{lstlisting}

\subsection{Комментарии}

Язык поддерживает однострочные и многострочные комментарии:

\begin{lstlisting}
// Однострочный комментарий

/* Многострочный
   комментарий */
\end{lstlisting}

\section{Пример программы}

Ниже приведен пример программы на языке \texttt{my\_lang}:

\begin{lstlisting}
class human {
    int id;
    string name;
    int hp;
    double money;
    
    void update() {
        hp = hp - 1;
    }
};

int main() {
    human p0;
    p0.id = 0;
    p0.name = "Alice";
    p0.hp = 100;
    p0.money = 50.0;
    
    int day = 0;
    while (day < 10) {
        p0.update();
        print("Day:", day, "HP:", p0.hp);
        day = day + 1;
    }
    
    return 0;
}
\end{lstlisting}

\section{Обработка ошибок}

Компилятор выводит сообщения об ошибках в следующем формате:

\begin{itemize}
    \item Синтаксические ошибки: указывается строка и позиция, где обнаружена ошибка
    \item Семантические ошибки: указывается тип ошибки и контекст
\end{itemize}

Пример сообщения об ошибке:

\begin{verbatim}
❌ Parsing failed (AST root is null).
❌ Semantic checks failed. Stopping compilation.
\end{verbatim}

\section{Ограничения}

В текущей версии компилятора существуют следующие ограничения:

\begin{itemize}
    \item Генерация POLIZ выполняется, но интерпретатор для выполнения POLIZ не реализован
    \item Некоторые встроенные функции могут быть не полностью реализованы
    \item Оптимизация AST ограничена простыми случаями
\end{itemize}

\section{Тестирование}

Для запуска тестов используйте команду:

\begin{lstlisting}[language=bash]
ctest
\end{lstlisting}

Или напрямую:

\begin{lstlisting}[language=bash]
./my_lang_tests.exe
\end{lstlisting}

\section{Часто задаваемые вопросы}

\subsection{Как добавить новое ключевое слово?}

Добавьте ключевое слово в файл \texttt{system\_reserved\_identifiers.txt} и пересоберите проект.

\subsection{Как изменить грамматику языка?}

Отредактируйте файл \texttt{grammar.txt} и пересоберите проект. Парсер будет автоматически перегенерирован.

\subsection{Где хранятся сгенерированные файлы?}

Все сгенерированные файлы (парсер, лексер) находятся в директории \texttt{out/}.

\section{Заключение}

Данное руководство описывает основные возможности компилятора языка \texttt{my\_lang}. Для получения дополнительной информации о внутреннем устройстве компилятора обратитесь к руководству программиста.

\end{document}

